% !TeX root = closed-systems-from-comparison-completeness.tex
% ============================================================
\chapter{Refinement Towers and the Quantifier-Reversal Boundary}
\label{ch:foundation}
% ============================================================
\section{Introduction}
Under the standing principle of closed-world admissibility
(\cref{sp:closed-world}), this chapter develops the compactness clause
\cref{sp:compactness} in explicit Boolean form.
It isolates the mechanism governing the passage from finite-stage
admissibility to global admissibility.
The Boolean-algebra and ultrafilter representation conventions used
below are the standard Stone conventions \cite{Stone1936}.
Its role in the sequel is foundational: it provides the exact
obstruction principle used later to pass from finite comparison data to
global comparison structure, and in particular underlies the Boolean
rectangle-splitting arguments of \cref{ch:bottom}.
Consider an increasing tower of finite-stage Boolean descriptions of a
system together with compatible exclusion ideals recording forbidden
local patterns. Because each stage ideal is proper, every finite stage
admits admissible ultrafilters, that is, ultrafilters avoiding the
excluded sets. The basic question is whether these finite-stage
admissible objects assemble coherently into a genuine admissible object
at the limit.
The obstruction is exact and is characterized by \cref{thm:boundary}.
For an exclusion tower, each stage ideal is proper, so
\cref{lem:ideal-sep} guarantees admissible ultrafilters at every
finite stage. The genuine question is whether these stagewise
admissible ultrafilters assemble into a global admissible ultrafilter
on \(\B_\infty\). This happens if and only if no finite family of
excluded sets from the tower already covers the whole universe.

Equivalently, under the standing exclusion-tower hypotheses,
\[
		\UF(\B_\infty;\J_\infty)\neq\varnothing
		\quad\Longleftrightarrow\quad
		\text{there do not exist }E_1,\dots,E_n\in\bigcup_{B\in\mathbb N}\J_B
		\text{ with }U=\bigcup_{i=1}^n E_i.
\]
Everything in this chapter is purely Boolean-algebraic and realizes only
the compactness layer of the standing principle.
No product structure, symmetry group, quotient semantics, transport
data, or smooth realization is invoked here; those belong to later
clauses \cref{sp:closure,sp:quotient,sp:transport}.
\medskip
Throughout this chapter,
\[
	\mathbb N=\{0,1,2,\dots\}.
\]
\section{Refinement towers}
\begin{definition}[Refinement tower]
	A \emph{refinement tower} consists of
	\begin{enumerate}[label=\textup{(\arabic*)}]
		\item a nonempty set \(U\);
		\item for each \(B\in\mathbb N\), a Boolean algebra
		      \[
			      \B_B\subseteq \mathcal P(U),
		      \]
		      viewed as a Boolean subalgebra of \(\mathcal P(U)\) with unit \(U\);
		\item inclusions
		      \[
			      \B_B\subseteq \B_{B+1}
			      \qquad
			      \text{for all }B\in\mathbb N.
		      \]
	\end{enumerate}
\end{definition}
We write
\[
	\B_\infty:=\bigcup_{B\in\mathbb N}\B_B.
\]
\begin{lemma}[The direct union is already Boolean]
	For a refinement tower, the union
	\[
		\B_\infty=\bigcup_{B\in\mathbb N}\B_B
	\]
	is a Boolean subalgebra of \(\mathcal P(U)\).
\end{lemma}
\begin{proof}
	Let \(A,C\in\B_\infty\). Choose \(B_A,B_C\in\mathbb N\) such that
	\[
		A\in\B_{B_A},
		\qquad
		C\in\B_{B_C}.
	\]
	Set
	\[
		B:=\max\{B_A,B_C\}.
	\]
	Since the tower is increasing,
	\[
		A,C\in\B_B.
	\]
	Because \(\B_B\) is a Boolean algebra, one has
	\[
		A\cup C\in\B_B,
		\qquad
		A\cap C\in\B_B,
		\qquad
		U\setminus A\in\B_B.
	\]
	Hence all three sets belong to \(\B_\infty\). Therefore \(\B_\infty\) is
	closed under finite unions, finite intersections, and complements, and is
	thus a Boolean subalgebra of \(\mathcal P(U)\).
\end{proof}
\begin{remark}[Direct-limit viewpoint]
	Because the tower is increasing, the direct limit of the Boolean algebras
	\((\B_B)_B\) is literally the union \(\B_\infty\). No additional Boolean
	completion is required.
\end{remark}
\section{Exclusion towers}
\begin{definition}[Exclusion tower]
	\label{def:excl-tower}
	Let
	\[
		\bigl(U,(\B_B)_{B\in\mathbb N}\bigr)
	\]
	be a refinement tower. An \emph{exclusion tower} over it is a family of
	ideals
	\[
		\J_B\subseteq \B_B
		\qquad
		(B\in\mathbb N)
	\]
	such that
	\begin{enumerate}[label=\textup{(\arabic*)}]
		\item each \(\J_B\) is proper;
		\item the family is compatible under restriction:
		      \[
			      \J_B=\J_{B+1}\cap \B_B
			      \qquad
			      \text{for all }B\in\mathbb N.
		      \]
	\end{enumerate}
\end{definition}
\begin{lemma}[Iterated compatibility]
	\label{lem:iter-compat}
	If \(B\le C\), then
	\[
		\J_B=\J_C\cap \B_B.
	\]
\end{lemma}
\begin{proof}
	The case \(C=B+1\) is part of \cref{def:excl-tower}. The general case
	follows by induction on \(C-B\).
\end{proof}
\begin{definition}[Limit ideal]
	The associated limit ideal is
	\[
		\J_\infty
		:=
		\left\langle
		\bigcup_{B\in\mathbb N}\J_B
		\right\rangle_{\B_\infty},
	\]
	that is, the ideal of \(\B_\infty\) generated by the union of the stage
	ideals.
\end{definition}
\begin{lemma}[Concrete description of the limit ideal]
	\label{lem:Jinf}
	Let
	\[
		\mathcal S:=\bigcup_{B\in\mathbb N}\J_B.
	\]
	Then
	\[
		\J_\infty
		=
		\left\{
		A\in\B_\infty:
		\exists n\ge1,\ \exists E_1,\dots,E_n\in\mathcal S
		\text{ such that }
		A\subseteq \bigcup_{i=1}^n E_i
		\right\}.
	\]
	In particular,
	\[
		U\in\J_\infty
		\quad\Longleftrightarrow\quad
		\exists n\ge1,\ \exists E_1,\dots,E_n\in\mathcal S
		\text{ such that }
		U=\bigcup_{i=1}^n E_i.
	\]
\end{lemma}
\begin{proof}
	Define
	\[
		\mathcal I
		:=
		\left\{
		A\in\B_\infty:
		\exists n\ge1,\ \exists E_1,\dots,E_n\in\mathcal S
		\text{ such that }
		A\subseteq \bigcup_{i=1}^n E_i
		\right\}.
	\]
	We show that \(\mathcal I\) is exactly the ideal generated by
	\(\mathcal S\).
	First, \(\mathcal I\) is an ideal of \(\B_\infty\). Certainly
	\[
		\varnothing\in\mathcal I,
	\]
	since \(\varnothing\subseteq E\) for every \(E\in\mathcal S\). If
	\(A\in\mathcal I\) and \(A'\in\B_\infty\) satisfies \(A'\subseteq A\),
	then the same witnessing family shows \(A'\in\mathcal I\). Thus
	\(\mathcal I\) is downward closed. If \(A,C\in\mathcal I\), choose
	witnesses
	\[
		A\subseteq \bigcup_{i=1}^m E_i,
		\qquad
		C\subseteq \bigcup_{j=1}^n F_j,
	\]
	with \(E_i,F_j\in\mathcal S\). Then
	\[
		A\cup C
		\subseteq
		\left(\bigcup_{i=1}^m E_i\right)\cup
		\left(\bigcup_{j=1}^n F_j\right),
	\]
	so \(A\cup C\in\mathcal I\). Hence \(\mathcal I\) is an ideal.
	Next, \(\mathcal S\subseteq\mathcal I\), since every \(S\in\mathcal S\)
	satisfies \(S\subseteq S\).
	Now let \(\mathcal K\) be any ideal of \(\B_\infty\) containing
	\(\mathcal S\). If \(A\in\mathcal I\), choose \(E_1,\dots,E_n\in\mathcal S\)
	such that
	\[
		A\subseteq \bigcup_{i=1}^n E_i.
	\]
	Since \(\mathcal K\) is an ideal containing each \(E_i\), it contains
	their finite union and hence, by downward closure, also \(A\). Therefore
	\(\mathcal I\subseteq\mathcal K\).
	Thus \(\mathcal I\) is the smallest ideal of \(\B_\infty\) containing
	\(\mathcal S\), so \(\mathcal I=\J_\infty\).
	The final equivalence follows by taking \(A=U\). Since each
	\(E_i\subseteq U\), the condition
	\[
		U\subseteq \bigcup_{i=1}^n E_i
	\]
	is equivalent to
	\[
		U=\bigcup_{i=1}^n E_i.
	\]
\end{proof}
\section{Ultrafilters and admissibility}
\begin{definition}[Ultrafilters and admissibility]
	For a Boolean algebra \(\B\), let \(\UF(\B)\) denote the set of
	ultrafilters on \(\B\). If \(\J\subseteq\B\) is an ideal, define
	\[
		\UF(\B;\J)
		:=
		\{\mathfrak u\in\UF(\B):\mathfrak u\cap\J=\varnothing\}.
	\]
	Elements of \(\UF(\B;\J)\) are called \emph{admissible ultrafilters}
	relative to \(\J\).
\end{definition}
\begin{definition}[Finite-stage and global admissibility]
	Finite-stage admissibility at level \(B\) means
	\[
		\UF(\B_B;\J_B)\neq\varnothing.
	\]
	Global admissibility means
	\[
		\UF(\B_\infty;\J_\infty)\neq\varnothing.
	\]
\end{definition}
\begin{lemma}[Ideal separation]
	\label{lem:ideal-sep}
	Let \(\B\) be a Boolean algebra and \(\J\subseteq\B\) an ideal. Then
	\[
		\UF(\B;\J)\neq\varnothing
		\quad\Longleftrightarrow\quad
		\J\text{ is proper}.
	\]
	Equivalently,
	\[
		\UF(\B;\J)=\varnothing
		\quad\Longleftrightarrow\quad
		1_\B\in\J.
	\]
\end{lemma}
\begin{proof}
	Assume first that \(\mathfrak u\in\UF(\B;\J)\). Since \(\mathfrak u\) is
	an ultrafilter,
	\[
		1_\B\in\mathfrak u.
	\]
	If \(1_\B\in\J\), then
	\[
		1_\B\in\mathfrak u\cap\J,
	\]
	contradicting admissibility. Hence \(\J\) is proper.
	Conversely, assume \(\J\) is proper, so that
	\[
		1_\B\notin\J.
	\]
	Define
	\[
		\mathcal F:=\{\neg E:E\in\J\}.
	\]
	We claim that \(\mathcal F\) has the finite intersection property. Let
	\(E_1,\dots,E_n\in\J\). Since \(\J\) is an ideal,
	\[
		E:=\bigvee_{i=1}^n E_i\in\J.
	\]
	Because \(\J\) is proper,
	\[
		E\neq 1_\B,
	\]
	and hence
	\[
		\neg E\neq 0_\B.
	\]
	By De Morgan's law,
	\[
		\bigwedge_{i=1}^n \neg E_i
		=
		\neg\left(\bigvee_{i=1}^n E_i\right)
		=
		\neg E
		\neq 0_\B.
	\]
	So \(\mathcal F\) has the finite intersection property.
	Let \(\mathfrak f\) be the filter generated by \(\mathcal F\). It is
	proper. By the ultrafilter extension theorem, there exists an ultrafilter
	\[
		\mathfrak u\supseteq\mathfrak f.
	\]
	If \(E\in\J\), then \(\neg E\in\mathcal F\subseteq\mathfrak u\). Thus
	\(E\notin\mathfrak u\), since otherwise
	\[
		E\wedge\neg E=0_\B\in\mathfrak u,
	\]
	which is impossible for a proper filter. Therefore
	\[
		\mathfrak u\cap\J=\varnothing,
	\]
	so \(\mathfrak u\in\UF(\B;\J)\).
\end{proof}
\section{Finite witnesses, first failure, and persistence}
\begin{lemma}[Finite witness]
	\label{lem:finite}
	Global admissibility fails if and only if there exist
	\[
		E_1,\dots,E_n\in\bigcup_{B\in\mathbb N}\J_B
	\]
	such that
	\[
		U=\bigcup_{i=1}^n E_i.
	\]
	Equivalently,
	\[
		\UF(\B_\infty;\J_\infty)=\varnothing
		\quad\Longleftrightarrow\quad
		\exists n\ge1,\ \exists E_1,\dots,E_n\in\bigcup_B\J_B
		\text{ with }
		U=\bigcup_{i=1}^n E_i.
	\]
\end{lemma}
\begin{proof}
	By \cref{lem:ideal-sep},
	\[
		\UF(\B_\infty;\J_\infty)=\varnothing
		\quad\Longleftrightarrow\quad
		\J_\infty\text{ is improper}.
	\]
	Since \(\J_\infty\) is an ideal in \(\B_\infty\), this is equivalent to
	\[
		U=1_{\B_\infty}\in\J_\infty.
	\]
	By \cref{lem:Jinf}, this holds if and only if there exist
	\(E_1,\dots,E_n\in\bigcup_B\J_B\) such that
	\[
		U\subseteq\bigcup_{i=1}^n E_i.
	\]
	Because each \(E_i\subseteq U\), this is equivalent to
	\[
		U=\bigcup_{i=1}^n E_i.
	\]
\end{proof}
\begin{remark}
	\cref{lem:finite} shows that global admissibility can fail even when
	each stage ideal \(\J_B\) is proper. The obstruction is not local but
	combinatorial: it arises exactly when finitely many excluded patterns,
	possibly drawn from different stages, already cover the entire universe.
\end{remark}
\begin{lemma}[First failure scale]
	\label{lem:firstscale}
	If global admissibility fails, then there exists \(B_0\in\mathbb N\) and
	sets
	\[
		E_1,\dots,E_n\in\J_{B_0}
	\]
	such that
	\[
		U=\bigcup_{i=1}^n E_i.
	\]
\end{lemma}
\begin{proof}
	By \cref{lem:finite}, choose
	\[
		E_1,\dots,E_n\in\bigcup_B\J_B
	\]
	such that
	\[
		U=\bigcup_{i=1}^n E_i.
	\]
	For each \(i\), choose \(B_i\in\mathbb N\) with \(E_i\in\J_{B_i}\), and
	set
	\[
		B_0:=\max_i B_i.
	\]
	Since the tower is increasing,
	\[
		E_i\in\B_{B_0}
		\qquad
		\text{for all }i.
	\]
	By \cref{lem:iter-compat},
	\[
		\J_{B_i}=\J_{B_0}\cap\B_{B_i}
		\qquad
		\text{for all }i,
	\]
	hence
	\[
		E_i\in\J_{B_0}
		\qquad
		\text{for all }i.
	\]
	Thus the same finite covering witness already occurs at stage \(B_0\).
\end{proof}
\begin{theorem}[Existence of a minimal failure depth]
	If global admissibility fails, then the set
	\[
		\mathcal B_{\mathrm{fail}}
		:=
		\left\{
		B\in\mathbb N:
		\exists E_1,\dots,E_n\in\J_B
		\text{ such that }
		U=\bigcup_{i=1}^n E_i
		\right\}
	\]
	is nonempty and therefore has a least element.
\end{theorem}
\begin{proof}
	By \cref{lem:firstscale}, the set \(\mathcal B_{\mathrm{fail}}\) is
	nonempty. Since it is a nonempty subset of \(\mathbb N\), it has a least
	element.
\end{proof}
\begin{theorem}[Persistence of failure]
	If admissibility fails at stage \(B_0\), then
	\[
		\UF(\B_B;\J_B)=\varnothing
		\qquad
		\text{for all }B\ge B_0.
	\]
\end{theorem}
\begin{proof}
	Assume there exist
	\[
		E_1,\dots,E_n\in\J_{B_0}
	\]
	such that
	\[
		U=\bigcup_{i=1}^n E_i.
	\]
	Fix \(B\ge B_0\). Because the tower is increasing,
	\[
		E_i\in\B_B
		\qquad
		\text{for all }i.
	\]
	By \cref{lem:iter-compat},
	\[
		\J_{B_0}=\J_B\cap\B_{B_0},
	\]
	hence
	\[
		E_i\in\J_B
		\qquad
		\text{for all }i.
	\]
	Therefore
	\[
		U=\bigcup_{i=1}^n E_i\in\J_B,
	\]
	so \(\J_B\) is improper. By \cref{lem:ideal-sep},
	\[
		\UF(\B_B;\J_B)=\varnothing.
	\]
\end{proof}
\section{The quantifier-reversal boundary}
\begin{definition}[Coherent ultrafilter tower]
	A \emph{coherent ultrafilter tower} is a sequence
	\[
		(\mathfrak u_B)_{B\in\mathbb N}
	\]
	such that
	\[
		\mathfrak u_B\in\UF(\B_B)
		\qquad
		\text{for all }B,
	\]
	and
	\[
		\mathfrak u_{B+1}\cap\B_B=\mathfrak u_B
		\qquad
		\text{for all }B.
	\]
	It is \emph{admissible} if
	\[
		\mathfrak u_B\cap\J_B=\varnothing
		\qquad
		\text{for all }B.
	\]
\end{definition}
\begin{lemma}[Union of a coherent ultrafilter tower]
	\label{lem:union-ultrafilter}
	Let
	\[
		(\mathfrak u_B)_{B\in\mathbb N}
	\]
	be a coherent ultrafilter tower, and define
	\[
		\mathfrak u:=\bigcup_{B\in\mathbb N}\mathfrak u_B\subseteq \B_\infty.
	\]
	Then \(\mathfrak u\) is an ultrafilter on \(\B_\infty\), and for every
	\(B\in\mathbb N\),
	\[
		\mathfrak u\cap\B_B=\mathfrak u_B.
	\]
\end{lemma}
\begin{proof}
	We first show that \(\mathfrak u\) is a proper filter on \(\B_\infty\).
	If \(A,C\in\mathfrak u\), choose \(B\) large enough that
	\(A,C\in\B_B\). Then \(A,C\in\mathfrak u_B\), hence
	\[
		A\cap C\in\mathfrak u_B\subseteq\mathfrak u.
	\]
	If \(A\in\mathfrak u\) and \(A\subseteq D\in\B_\infty\), choose \(B\)
	large enough that \(A,D\in\B_B\). Then \(A\in\mathfrak u_B\), so by
	upward closure of the filter \(\mathfrak u_B\),
	\[
		D\in\mathfrak u_B\subseteq\mathfrak u.
	\]
	Also \(\varnothing\notin\mathfrak u\), since \(\varnothing\notin\mathfrak u_B\)
	for every \(B\). Thus \(\mathfrak u\) is a proper filter on
	\(\B_\infty\).
	Next we show that \(\mathfrak u\) decides every element of \(\B_\infty\).
	Let \(A\in\B_\infty\). Choose \(B\) such that \(A\in\B_B\). Since
	\(\mathfrak u_B\) is an ultrafilter on \(\B_B\), exactly one of \(A\)
	and \(U\setminus A\) belongs to \(\mathfrak u_B\), hence exactly one
	belongs to \(\mathfrak u\). Therefore \(\mathfrak u\) is an ultrafilter
	on \(\B_\infty\).
	Finally, if \(A\in \mathfrak u\cap \B_B\), then \(A\in\mathfrak u_C\) for
	some \(C\ge B\). Coherence gives
	\[
		\mathfrak u_C\cap\B_B=\mathfrak u_B,
	\]
	so \(A\in\mathfrak u_B\). The reverse inclusion
	\[
		\mathfrak u_B\subseteq \mathfrak u\cap\B_B
	\]
	is immediate. Hence
	\[
		\mathfrak u\cap\B_B=\mathfrak u_B
		\qquad
		\text{for all }B.
	\]
\end{proof}
\begin{lemma}[Admissibility passes to the limit]
	\label{lem:admiss-limit}
	Let
	\[
		(\mathfrak u_B)_{B\in\mathbb N}
	\]
	be an admissible coherent ultrafilter tower, and let
	\[
		\mathfrak u:=\bigcup_{B\in\mathbb N}\mathfrak u_B.
	\]
	Then
	\[
		\mathfrak u\in\UF(\B_\infty;\J_\infty).
	\]
\end{lemma}
\begin{proof}
	By \cref{lem:union-ultrafilter}, \(\mathfrak u\) is an ultrafilter on
	\(\B_\infty\). It remains to prove admissibility.
	Let \(A\in\J_\infty\). By \cref{lem:Jinf}, there exist
	\[
		E_1,\dots,E_n\in\bigcup_B\J_B
	\]
	such that
	\[
		A\subseteq\bigcup_{i=1}^n E_i.
	\]
	Choose \(B\) with all \(E_i\in\J_B\). Since
	\[
		\mathfrak u\cap\B_B=\mathfrak u_B
	\]
	by \cref{lem:union-ultrafilter}, and
	\[
		\mathfrak u_B\cap\J_B=\varnothing,
	\]
	each \(E_i\notin\mathfrak u\), hence
	\[
		U\setminus E_i\in\mathfrak u
		\qquad
		\text{for all }i.
	\]
	Therefore
	\[
		\bigcap_{i=1}^n(U\setminus E_i)\in\mathfrak u.
	\]
	But
	\[
		\bigcap_{i=1}^n(U\setminus E_i)\subseteq U\setminus A,
	\]
	so
	\[
		U\setminus A\in\mathfrak u,
	\]
	and hence \(A\notin\mathfrak u\). Thus
	\[
		\mathfrak u\cap\J_\infty=\varnothing,
	\]
	so
	\[
		\mathfrak u\in\UF(\B_\infty;\J_\infty).
	\]
\end{proof}
\begin{theorem}[Quantifier-reversal boundary]
	\label{thm:boundary}
	The following are equivalent.
	\begin{enumerate}[label=\textup{(\arabic*)}]
		\item
		      \[
			      \UF(\B_\infty;\J_\infty)\neq\varnothing.
		      \]
		\item For every finite family
		      \[
			      E_1,\dots,E_n\in\bigcup_{B\in\mathbb N}\J_B,
		      \]
		      one has
		      \[
			      \bigcup_{i=1}^n E_i\neq U.
		      \]
		\item There exists an admissible coherent ultrafilter tower.
	\end{enumerate}
\end{theorem}
\begin{proof}
	\emph{(1)\(\Rightarrow\)(2).}
	This is the contrapositive of \cref{lem:finite}.
	\smallskip
	\emph{(2)\(\Rightarrow\)(3).}
	We construct an admissible coherent tower inductively.
	For the base step, condition \textup{(2)} implies that \(\J_0\) is
	proper; otherwise
	\[
		U\in\J_0\subseteq\bigcup_B\J_B,
	\]
	contradicting \textup{(2)}. By \cref{lem:ideal-sep}, choose
	\[
		\mathfrak u_0\in\UF(\B_0;\J_0).
	\]
	Assume \(\mathfrak u_B\in\UF(\B_B;\J_B)\) has been constructed. Define
	\[
		\mathcal G_{B+1}
		:=
		\mathfrak u_B\cup\{U\setminus E:E\in\J_{B+1}\}
		\subseteq \B_{B+1}.
	\]
	We claim that \(\mathcal G_{B+1}\) has the finite intersection property.
	Take finitely many
	\[
		A_1,\dots,A_m\in\mathfrak u_B,
		\qquad
		E_1,\dots,E_n\in\J_{B+1},
	\]
	and set
	\[
		A:=\bigcap_{j=1}^m A_j\in\mathfrak u_B,
		\qquad
		E:=\bigcup_{i=1}^n E_i\in\J_{B+1}.
	\]
	If
	\[
		A\cap(U\setminus E)=\varnothing,
	\]
	then \(A\subseteq E\). Since \(A\in\B_B\subseteq\B_{B+1}\) and
	\(\J_{B+1}\) is downward closed, this implies
	\[
		A\in\J_{B+1}\cap\B_B=\J_B,
	\]
	contradicting \(\mathfrak u_B\cap\J_B=\varnothing\). Therefore
	\[
		A\cap(U\setminus E)\neq\varnothing.
	\]
	So \(\mathcal G_{B+1}\) has the finite intersection property.
	Let \(\mathfrak f_{B+1}\) be the filter generated by
	\(\mathcal G_{B+1}\). It is proper, hence extends to an ultrafilter
	\[
		\mathfrak u_{B+1}\in\UF(\B_{B+1}).
	\]
	Since \(\mathfrak u_B\subseteq\mathfrak u_{B+1}\) and
	\(\mathfrak u_B\) is already an ultrafilter on \(\B_B\), one has
	\[
		\mathfrak u_{B+1}\cap\B_B=\mathfrak u_B.
	\]
	Also, for every \(E\in\J_{B+1}\),
	\[
		U\setminus E\in\mathfrak u_{B+1},
	\]
	so \(E\notin\mathfrak u_{B+1}\). Hence
	\[
		\mathfrak u_{B+1}\cap\J_{B+1}=\varnothing.
	\]
	Induction yields an admissible coherent ultrafilter tower.
	\smallskip
	\emph{(3)\(\Rightarrow\)(1).}
	This is \cref{lem:admiss-limit}.
\end{proof}
\section{Inverse-limit representation}
\begin{theorem}[Inverse-limit representation]
	\label{thm:inverse-limit}
	Restriction along the inclusions \(\B_B\subseteq\B_{B+1}\) induces a
	canonical bijection
	\[
		\UF(\B_\infty)\cong \varprojlim_B \UF(\B_B),
	\]
	where the inverse limit is taken with respect to the restriction maps
	\[
		\UF(\B_{B+1})\to\UF(\B_B).
	\]
\end{theorem}
\begin{proof}
	Define
	\[
		\Phi:\UF(\B_\infty)\to \varprojlim_B \UF(\B_B),
		\qquad
		\Phi(\mathfrak u):=(\mathfrak u\cap\B_B)_B.
	\]
	This is well defined because the restriction of an ultrafilter to a
	Boolean subalgebra is an ultrafilter, and the restrictions are coherent.
	Conversely, let
	\[
		(\mathfrak u_B)_B\in\varprojlim_B \UF(\B_B).
	\]
	Define
	\[
		\Psi((\mathfrak u_B)_B):=\bigcup_B\mathfrak u_B\subseteq\B_\infty.
	\]
	By \cref{lem:union-ultrafilter}, this union is an ultrafilter on
	\(\B_\infty\), and for every \(B\),
	\[
		\Psi((\mathfrak u_B)_B)\cap\B_B=\mathfrak u_B.
	\]
	Thus \(\Psi\) is a well-defined map
	\[
		\Psi:\varprojlim_B \UF(\B_B)\to \UF(\B_\infty).
	\]
	It is immediate that
	\[
		\Phi(\Psi((\mathfrak u_B)_B))=(\mathfrak u_B)_B
	\]
	and
	\[
		\Psi(\Phi(\mathfrak u))=\mathfrak u.
	\]
	Therefore \(\Phi\) is a bijection with inverse \(\Psi\).
\end{proof}
\begin{corollary}[Admissible inverse-limit representation]
	\label{cor:admiss-inv}
	Restriction induces a canonical bijection
	\[
		\UF(\B_\infty;\J_\infty)
		\cong
		\varprojlim_B \UF(\B_B;\J_B).
	\]
\end{corollary}
\begin{proof}
	By \cref{thm:inverse-limit}, ultrafilters on \(\B_\infty\) are in
	bijection with coherent ultrafilter towers on the stages.
	If
	\[
		\mathfrak u\in\UF(\B_\infty;\J_\infty),
	\]
	then for each \(B\),
	\[
		\mathfrak u\cap\J_B=\varnothing,
	\]
	because \(\J_B\subseteq\J_\infty\). Hence each stage restriction is
	admissible.
	Conversely, let \((\mathfrak u_B)_B\) be a coherent tower with
	\[
		\mathfrak u_B\in\UF(\B_B;\J_B)
		\qquad
		\text{for all }B.
	\]
	Then by \cref{lem:admiss-limit},
	\[
		\bigcup_B\mathfrak u_B\in\UF(\B_\infty;\J_\infty).
	\]
	Thus the bijection of \cref{thm:inverse-limit} restricts to the claimed
	bijection on admissible ultrafilters.
\end{proof}
\section{Stability criterion}
\begin{theorem}[Stability of admissibility]
	\label{thm:stability}
	For an exclusion tower satisfying the compactness clause
	\cref{sp:compactness}, finite-stage admissibility holds at every level,
	and global admissibility holds if and only if finite non-coverage
	holds, namely condition \textup{(2)} of \cref{thm:boundary}.
\end{theorem}
\begin{proof}
	Because each stage ideal \(\J_B\) is proper by
	\cref{def:excl-tower}, \cref{lem:ideal-sep} gives
	\(\UF(\B_B;\J_B)\neq\varnothing\) for every \(B\).
	The remaining claim is exactly \cref{thm:boundary}, which identifies
	global admissibility with finite non-coverage without any further
	hypothesis.
\end{proof}
\section{Conclusion}
\label{sec:foundation-conclusion}
The compactness boundary established in
\cref{lem:finite,lem:firstscale,thm:boundary,cor:admiss-inv} is exact:
for an exclusion tower, each proper stage ideal already admits
admissible ultrafilters, and global admissibility fails if and only if
finitely many exclusions from the tower already cover the universe.
Equivalently, admissible ultrafilters on the limit Boolean algebra are
precisely admissible coherent towers of stage ultrafilters, so once
finite non-coverage holds, no further global obstruction remains.

\Cref{ch:bottom} applies this compactness boundary to comparison worlds
and derives the canonical factorization that initiates quotient semantics.
